\subsection{Eliminación de las cuatro líneas de la casilla insertada}\label{sec:B4T01}
El padre es un residual ordinario con $k\ge\delta_0n$ y un punto estacionario
de $H_B$ con $0<p_e\le C/k$, vacancias diagonales $h\ge\eta$ y brecha del
laplaciano RC ponderado al menos $\gamma$. Las constantes son fijas.
Tomamos $k$ suficientemente grande para que $p<1/4$.
Sean $q=1-p$, $\tau=pq/(1-4p)$, $T=\operatorname{diag}\tau$ y
\[
 S=\operatorname{diag}(0_{RC},h)+A^TTA,\quad
 v=(\mathbf1_R,-\mathbf1_C,0_D),\quad G=S^+.
\]
De las Secciones~\ref{sec:B3T05}--\ref{sec:B3T07} usamos la brecha uniforme de $S$,
$G_{ij}=\delta_{ij}+O(k^{-1})$ y $\|G\|_\infty=O(1)$.
Para el gauge usamos un cofactor RC, no un pseudodeterminante sin normalizar.

\paragraph{Las cuatro líneas eliminadas y sus casillas.}
Sean $J$ las cuatro líneas de $c$, $I=J^c$, $F$ las casillas de su unión
y $E'=E\setminus F$. Las coordenadas siguen siendo las originales.
Cada línea retenida pierde a lo sumo tres casillas. Escribimos
\[
 d_\ell=\sum_{e\in F\cap\ell}p_e,\quad w=p|_{E'},\quad
 U=A_{F,I}^TT_FA_{F,I},\quad \Delta=\operatorname{diag}_{D\cap I}d,
\]
\[
 Q=S_{IJ}S_{JJ}^{-1}S_{JI},\quad P_c=S_{II}-Q,\quad E_c=Q+\Delta-U.
\]
Entonces valen las identidades exactas $S_{\rm hijo}(w)=P_c+E_c$ y
$P_cv_I=E_cv_I=0$. La matriz $Q$ tiene rango a lo sumo cuatro;
en cambio, $\Delta-U$ no tiene por qué tener rango acotado.

Para el punto factible $\bar p=w+z$ de la Sección~\ref{sec:B3T07} se añade
\[
 \Xi=A'^T\operatorname{diag}(\tau(\bar p)-\tau(w))A'
       -\operatorname{diag}_{D\cap I}(D'^Tz),\quad
 \mathcal E=E_c+\Xi,\quad S_{\rm hijo}(\bar p)=P_c+\mathcal E.
\]
Aquí $D'$ es la incidencia casilla--diagonal del hijo.
Para $y\perp v_I$, minimizar sobre las cuatro coordenadas eliminadas da
\[
 y^TP_cy=\min_a(a,y)^TS(a,y)\ge c_*\|y\|^2.
\]
En efecto, $\operatorname{dist}((a,y),\operatorname{span}v)\ge\|y\|$.
Por tanto, el núcleo de $P_c$ es exactamente $\operatorname{span}v_I$.

\paragraph{Determinante exacto.}
Escogemos la coordenada RC del gauge dentro de $I$ y definimos
\[
 W_c=P_c^{+1/2}\mathcal E P_c^{+1/2}\big|_{v_I^\perp}.
\]
El complemento de Schur y, después, la identidad de cofactores sobre un mismo núcleo dan
\[
 \boxed{\frac{\det S_{{\rm hijo},\rm gf}(\bar p)}{\det S_{{\rm padre},\rm gf}(p)}
 =\frac{\det_{v_I^\perp}(I+W_c)}{\det S_{JJ}}.}
\]
Si se usan pseudodeterminantes, el segundo miembro debe multiplicarse por $(k-1)/k$,
pues $\operatorname{pdet}S=2k\det S_{\rm gf}$ y el hijo tiene $k-1$ filas.
No se añade ese factor una segunda vez al trabajar con cofactores.

\paragraph{La misma eliminación en $N$.}
Con un gauge retenido, $N=\left(\begin{smallmatrix}T^{-1}&A_{\rm gf}\\
A_{\rm gf}^T&-\operatorname{diag}h_{\rm gf}\end{smallmatrix}\right)$.
El bloque eliminado $(F,J)$ y su inverso parcial son
\[
 K_N=\begin{pmatrix}T_F^{-1}&A_{F,J}\\A_{F,J}^T&-\operatorname{diag}h_J\end{pmatrix},
 \quad \det K_N=\Big(\prod_F\tau_e^{-1}\Big)\det S_{JJ},
\]
\[
 (K_N^{-1})_{FF}=T_F-T_FA_{F,J}S_{JJ}^{-1}A_{F,J}^TT_F.
\]
El complemento de Schur retenido tiene bloque inferior $-\operatorname{diag}h_I-U+Q$.
Para pasar a $N_{\rm hijo}(w)$ se resta $E_c$ de ese bloque inferior.
Para $\bar p$, se cambian además $T_{E'}^{-1}$ y las vacancias.
Se recupera $\mathcal D=\prod_e[(1-4p_e)/q_e^4]\det S_{\rm gf}$.
No se ha usado ningún desarrollo de Taylor que lleve una coordenada eliminada hasta cero o uno.

\paragraph{Resto matricial uniforme.}
Las incidencias geométricas y la Sección~\ref{sec:B3T07} dan
$\|\mathcal E\|_2=O(k^{-1})$ y $\|\mathcal E\|_F^2=O(k^{-1})$:
$\Delta-U$ es dispersa por filas; $Q$ tiene rango a lo sumo cuatro y norma $O(k^{-1})$;
$\Xi$ tiene suma absoluta por fila $O(k^{-1})$ y dimensión $O(k)$.
Así $\|W_c\|_2=O(k^{-1})$, $\operatorname{tr}W_c^2=O(k^{-1})$ y
\[
 \boxed{\log\det(I+W_c)=\operatorname{tr}W_c
       -\tfrac12\operatorname{tr}W_c^2+O(k^{-2}).}
\]
El resto se acota por $C\|W_c\|_2\operatorname{tr}W_c^2$.
El término cuadrático no puede descartarse como $O(k^{-2})$.
